% !TEX root = ../main.tex
\hspace{24pt}

\section{Conclusions}

This study has proposed TCP-Uni3C, a complete end-to-end pipeline that takes a single TCP image together with a small number of manual depth annotations as input
and generates a 2.5D video with convincing parallax,
without constructing 3D models or collecting large-scale TCP training datasets,
thereby enabling static traditional Chinese paintings to be presented in animated form.

To address the domain gap of existing depth estimation models on TCP,
this study adopts a three-stage human-in-the-loop strategy:
a discriminative model first produces a rough depth map for the user to inspect;
the user then annotates sparse depth control points at key positions;
finally, the control points guide the denoising optimization process of Marigold-DC to refine the depth field.
Within this process, we propose \textit{ChinesePaintingLoss}, a composite loss function designed specifically for the characteristics of TCP,
in which a brushstroke-aware smoothness loss eliminates depth discontinuities on both sides of dense-ink brushstrokes
and a whitespace regularization term pushes the white-space background toward the far distance,
so that the refined depth map simultaneously satisfies the numerical constraints of the control points and the spatial-semantic conventions of Chinese painting.
Ablation studies confirm that each loss component makes a clear and complementary contribution to depth quality.

After obtaining the refined depth map,
this study unprojects it into a three-dimensional point cloud,
derives the initial camera distance from the principle of similar triangles,
and designs a circular camera trajectory whose start and end are smoothly joined to render the point-cloud video,
which serves as the geometric guidance condition for Uni3C;
combined with text prompts designed for static TCP scenes,
the video diffusion model is driven to fill the holes exposed by viewpoint shifts while preserving the ink-and-brush style.
Experimental results show that,
compared with the original Uni3C driven by its built-in depth estimation,
the videos generated by TCP-Uni3C exhibit pronounced and stable foreground--background parallax,
convey a genuine sense of spatial depth,
and preserve the visual style of the original painting across a variety of subjects, including bird-and-flower, animal, and figure paintings.

\section{Future Work}

Building on the results of this study and the limitations discussed in Chapter 4,
future research may proceed in the following directions:

\begin{enumerate}
  \item \textbf{Constructing a multi-view TCP depth dataset}:
        The scarcity of TCP depth annotation data is the fundamental reason this study adopts a human-in-the-loop strategy,
        yet the proposed pipeline itself can serve as a dataset construction tool---
        for every painting corrected with manual control points,
        the pipeline simultaneously produces a refined depth map, a three-dimensional point cloud,
        and a multi-view image sequence rendered along the camera trajectory.
        By systematically accumulating such annotation results over a large number of paintings,
        a multi-view depth dataset for the TCP domain can be constructed,
        allowing subsequent research to train or fine-tune depth estimation models dedicated to TCP,
        gradually reducing the reliance on manual annotation,
        and providing a data foundation for further related research such as 3D reconstruction of TCP and stylized novel view synthesis.

  \item \textbf{Layered depth representation}:
        To resolve the depth ambiguity in regions of overlapping brushstrokes,
        layered depth image (LDI) or semi-transparent layer decomposition techniques could be introduced
        to decompose regions of blended ink into multiple depth layers,
        so that foreground and background objects exhibit correct occlusion and parallax relations under viewpoint shifts.

  \item \textbf{Reducing manual annotation cost}:
        The current pipeline still relies on the user to manually click control points.
        Future work could incorporate semantic segmentation or vision-language models
        to automatically identify the white space, main subjects, and depth layers in a painting and to recommend control-point positions and depth values,
        so that the user only needs to confirm or fine-tune the recommendations, further lowering the operational barrier.

  \item \textbf{Automatic determination of camera parameters}:
        As shown by the experiments in Chapter 4, parameters such as the camera distance scaling and the angular amplitude have a significant impact on generation quality,
        and they currently must be tuned manually according to the composition of each painting.
        Future work could train a camera parameter prediction model
        that takes the painting content and the geometric distribution of the point cloud as input
        and automatically estimates suitable settings such as the camera distance and the trajectory amplitude,
        so that the video generation process no longer requires manual parameter adjustment,
        further improving the degree of automation of the overall pipeline.

\end{enumerate}
